%% ---------------------------------------------------------------------
\section{Ablations and controls}\label{sec:ablations}

\looseness=-1 The benchmark's defensibility rests on knowing \emph{which} ingredient carries
the gain, so we ablate every one, moving outward from the ingredient itself to
the design choices around it. First the target, against random, learned and
alternative hand-crafted ones (Section~\ref{abl:target}), then where it is
applied, how its weight is scheduled and what form the head takes
(Section~\ref{abl:placement}). Next, whether one configuration survives changes
of architecture and depth without retuning (Section~\ref{abl:universal}), the
classical alternative of placing the same bank in the forward path and its
failure mode (Section~\ref{abl:forward}), and whether a wider bank helps
(Section~\ref{abl:width}). Last, a second effect the accuracy tables alone
would conflate with the first: the prior's stabilization of backbones whose
baselines collapse (Section~\ref{abl:stab}). Section~\ref{abl:notablated}
states what we did \emph{not} vary.

\subsection{The target is the ingredient}\label{abl:target}
\begin{table}[pos=htbp]
\caption{Target and mechanism controls on CIFAR-100 with ResNet-18 (each
row is the champion configuration with one ingredient varied; $\Delta$ vs.\
the shared baseline, three seeds in the varied arm against the released
ten-seed baseline). These cells hold $\lambda{=}0.3$, not the champion's decaying schedule, so they are comparable to each other, as an
ablation requires, but not to Table~\ref{tab:envelope}. The gain needs the
moment structure specifically: random fixed targets do nothing, a
$2\times$-cost learned teacher does nothing anywhere on a nine-point envelope,
and the classical HOG descriptor recovers a third to a half.}
\label{tab:controls}
\centering\scriptsize
\setlength{\tabcolsep}{2.6pt}
\zebra{4}
\begin{tabular}{P{0.33\columnwidth}rrP{0.31\columnwidth}}
\toprule
 & \multicolumn{2}{c}{$\Delta$ at fraction~\up} & \\
\cmidrule(lr){2-3}
Aux target & 5\% & 10\% & Note \\
\midrule
Moment magnitude (ours) & $+3.18$ & $+2.71$ & phase-invariant energy \\
Structure tensor & $+1.08$ & $+1.68$ & mild nonlinearity \\
Steerable harmonics & $+1.55$ & $+0.91$ & principled rot.-inv. \\
Rotation-invariant energy & $+1.44$ & $+0.74$ & \\
Oriented edges (Gabor) & $+0.56$ & $-0.09$ & raw edges do not help \\
2nd-order invariants & $-0.43$ & $-3.07$ & over-processed \\
\midrule
HOG target & $+1.08$ & $+1.27$ & descriptor helps; moments ${\sim}2.5\times$ better \\
Learned teacher (FitNets) & -- & -- & $2\times$ cost; $-0.29$ to $+0.73$ over nine fractions, no 5/10\% cell \\
Random fixed maps & $-0.12$ & $+0.04$ & kills ``any aux signal'' \\
\midrule
Cosine loss (not MSE) & $+0.62$ & $+0.36$ & magnitude scale matters \\
\bottomrule
\end{tabular}
\end{table}

\looseness=-1 Table~\ref{tab:controls}: a random fixed target of identical shape moves
accuracy by at most $0.12$ points in either direction, so the mechanism is not
``regression as regularization''. A FitNets teacher trained on the same data at twice the
cost does nothing anywhere we measured it (nine fractions from $1$ to
$100\%$, $-0.29$ to $+0.73$), making the free hand-crafted target strictly better than an
expensive learned one in this regime. Among hand-crafted families, phase-invariant \emph{magnitude} wins at both
fractions by a factor of $1.6$ to $2$ over the next-best target, and
discarding its magnitude scale (cosine loss) forfeits most of the gain. The
ordering \emph{below} magnitude is not stable, and we would rather say so than
rank six targets off one column: the structure tensor is fourth at $5\%$ and
second at $10\%$, and only the two ends hold at both fractions, magnitude
first and heavily processed invariants last and actively harmful.

\subsection{Placement, schedule, and head}\label{abl:placement}
The tap depth is flat across the first three stages (within seed noise at
both 1\% and 10\%) with a cliff only at the final stage: the one design
knob in the study that is \emph{not} a data-regime knob. The schedule
start $\lambda_0$ \emph{is} a regime knob (best $2.0$ at 1--2\%, $1.0$ at
3--10\%, $0.3$ at 15--25\%), but its reach is bounded: three attempts to
push $\Delta$ past the measured feature gain by retuning $\lambda_0$ all
failed. Consistent with the law, the knob cannot manufacture $G$; it can only
realize it. Head-norm projection, derived from a traced scale
degeneracy (the aux loss is invariant under feature-shrink with head-inflate,
which collapses bottleneck ResNets), converts ResNet-50 from a bistable
failure ($\{41.2, 36.4, 42.4\}$ across seeds) into a stable $+3.9$
($\{44.5, 45.0, 44.4\}$) and is free elsewhere; we adopt it always-on.
Three plausible fixes tried \emph{before} tracing the mechanism (gradient
balancing, disabling automatic mixed precision, BatchNorm on the tap) all
failed or backfired, the last making collapse $25\times$ faster: a small case
study in mechanism-first debugging.

\subsection{One configuration, every architecture and depth}\label{abl:universal}
A single champion setting transplants without retuning: at CIFAR-100 with
$10\%$ of the data, ResNet-18, -34 and -50 land $+4.3$, $+4.0$ and $+3.9$
against their own baselines (these are the head-norm cells of the
backbone sweep, not the frozen-recipe envelope of
Table~\ref{tab:envelope}); the new depth
sweep confirms one $\lambda_0$ across ResNet-34 and -50 on six further
domains ($+0.4$ to $+9.0$
at 5--10\% everywhere except CUB's deep-scarcity floor, decaying on the
right flank as the law requires). Head \emph{form} is irrelevant where it
plausibly might not have been: linear, cosine and nearest-class-mean readouts
measure the same aux-vs-baseline gap to within $0.3$ points. The left-flank
bottleneck is label information, not classifier expressivity.

\subsection{The forward-path alternative}\label{abl:forward}
Placing the same bank \emph{in the forward path} as concatenated fixed
channels at the stem, the classical hybrid design, yields the study's largest
structural contrast. Forward-path stems are the low-data
accuracy leaders (energy-magnitude reaches $+2.6$ and $+3.5$ at CIFAR-100
1--2\%, beating the auxiliary prior there), but every fixed forward
variant we test (linear Gabor at two kernel scales, and four nonlinear energy
families) pays a penalty band at 10--25\% data ($-1.4$ to $-7.3$) that no
filter choice escapes, and the phase-invariant variants
never wash out even at 100\% ($-3.9$): a hard constraint permanently
occupies input bandwidth that abundant data would rather use. The
auxiliary placement of the \emph{same} knowledge is what removes the
band, which is the design argument for fusing priors through
training-only targets rather than through architecture.

\subsection{Bank width}\label{abl:width}
Widening the pinned 8-pair bank to 12 channels (via a third octave or two
extra orientations, indistinguishable at $10{\times}10$ seeds) buys a
real but local $+0.3$--$0.5$ at Tiny-ImageNet's 64\,px and nothing at
32\,px; across five further 64\,px domains the pooled excess fails our
pre-registered re-pin criterion. The pinned 8-pair bank therefore
stays, with cross-domain evidence, not convenience, behind it.

\subsection{Stabilization as a second, distinct effect}\label{abl:stab}
Across ConvNeXt (under the frozen SGD recipe), ResNet-50 (without
head-norm), Swin-T (on five datasets), and Swin at ImageNet64, baselines
are bistable, some seeds train, some sit at chance, while the prior
arm trains tightly every time. Swin's feature-level variance shows the
instability lives in the \emph{features} (baseline evaluation $\sigma$ up to
$4.3$ vs.\ $0.2$ with the prior). We report stabilization as a real,
separate property; we do not claim routine variance reduction (a
$10{\times}10$-seed test of that claim came back negative, and we report
that too).

%% ---------------------------------------------------------------------
\subsection{What we did not ablate}\label{abl:notablated}
Completeness cuts both ways, so we list the design constants this paper does
\emph{not} vary. Four are applied identically to both arms of every pair, so they bound absolute
values without being able to explain a paired difference. Two concern training:
the recipe, which is the instrument and is never tuned
(Section~\ref{sec:recipe}), and average pooling of the target to the tapped
stage's resolution, with max and learned pooling untested. Two concern what
reads it: the auxiliary head, a single $1{\times}1$ convolution we never widen
or deepen, although we ablate the \emph{classifier} head thoroughly
(Section~\ref{abl:placement}), and the linear evaluation's ridge constant of
$10^{-4}$, which shifts every arm together, so $G$'s sign and ordering are
unaffected while its numeric value is conditional on it.

\looseness=-1 Three matter more, because they define the knowledge source, not the measurement. The bank's envelope width is tied to its frequency by a fixed
relation, so every filter carries the same number of cycles; we vary how many
scales and orientations the bank holds (Section~\ref{abl:width}) but never
that constant, leaving ``how localized should an oriented-energy target be''
open. The target is computed on BT.601 luminance, which makes the prior
colorblind by construction, and a color-opponent target is untested. And while
the schedule's endpoint is ablated and load-bearing ($\lambda(T)=0$ against
$\lambda(T)=0.1$), its cosine \emph{shape} is not. None of these is a confound
here, since a paired $\Delta$ and a same-protocol $G$ gap are differences
taken under identical settings, but they bound what one bank, one head form
and one schedule shape can establish, and the first two are the natural next
dimensions for anyone extending the taxonomy to other structural priors.

%% ---------------------------------------------------------------------
\section{A decision guide}\label{sec:guide}

\begin{table}[pos=htbp]
\caption{What to use when: the grid distilled to recommendations. Costs
are training-compute multiples of the baseline; gains are typical measured
ranges, not promises. ``Prior'' is the champion MomentAux configuration
verbatim.}
\label{tab:guide}
\centering\scriptsize
\setlength{\tabcolsep}{2.6pt}
\zebra{2}
\begin{tabular}{P{0.285\columnwidth}P{0.205\columnwidth}lP{0.265\columnwidth}}
\toprule
Regime & Recommendation & Cost & Measured basis \\
\midrule
Conv backbone, photo-like data, mid band (2--15\%), $2\times$ affordable & SimCLR init & $2\times$ & $+2..{+}9$ over baseline; beats prior by $+2..{+}5$ \\
Conv, photo-like, $2\times$ \emph{not} affordable & prior & $1.02\times$ & $+1.4..{+}6.7$; ${\sim}$SSL parity at extremes \\
Conv, extreme scarcity (${\leq}1$--$2\%$) & prior (or fwd stem) & $1.02\times$ & SSL margin ${\leq}1$ pt; forward stem best ${\leq}2\%$ \\
Satellite-like statistics, low data & prior & $1.02\times$ & beats SimCLR at 7 of 11 EuroSAT fractions \\
Any ViT trained the modern way (DeiT aug) & prior \emph{and} augmentation & $1.02\times$ & flip: beats SSL at every fraction at $2\times$; at $5\times$ only in the higher-data band (both populations); $+14..{+}25$ over recipe alone \\
ViT and Swin from scratch, every scale measured & prior & $1.02\times$ & $+3.2$ at 1.28M imgs; $+13$ and $+26$ at 224\,px; stabilizes Swin \\
ImageNet-transferable task & transfer, without the prior &--& tax up to $-17$ under the tested schedule, feature-side \\
Domain-shifted from ImageNet (e.g.\ stains) & prior from scratch competitive & $1.02\times$ & transfer advantage shrinks to $+2..{+}3$ \\
Conv, photo-like, ${\geq}50\%$ & nothing (all converge) &--& conv arms within $\pm 0.8$; \emph{not} ViTs, which still gain at full data \\
\bottomrule
\end{tabular}
\end{table}

Table~\ref{tab:guide} is the benchmark's practical distillate. Its
strongest single row is the modern-ViT one: for anyone training a small
or mid-sized ViT with the standard recipe, a fixed spectral auxiliary
target at $2\%$ overhead is, in every regime we measured, either the best
or tied-best intervention available, and it composes with the recipe
rather than competing with it.

%% ---------------------------------------------------------------------
\section{Discussion}\label{sec:discussion}

\paragraph{What the decomposition buys, and what it does not.}
\looseness=-1 The grid's size is not its point; its point is that one low-dimensional
structure organizes much of it. We are careful with ``explains'': as
Section~\ref{sec:lawform} quantifies, this is an account of a trend and of a
\emph{sign}, not a predictor of a cell's value. Eq.~\ref{eq:law} says an intervention's
visible benefit is feature gain passed through a readout channel set by how
well the task is already being solved. One mechanism then covers observations
that otherwise need separate stories: why gains collapse at extreme
scarcity even when features improve (readout deeply negative, five
labels per class cannot express a better representation); why relabeling
the same pixels with coarser classes triples the visible gain (the
baseline crosses the readout zero, not because features changed; we
measure them unchanged); why every intervention converges at sufficiency
(both terms go to zero); and why coarse label spaces read ${\sim}2$ points more gain
under coarse evaluation than fine on the same checkpoints (the measuring
stick itself is label-hungry).

\paragraph{What the currency account contributes to fusion theory.}
\looseness=-1 None of our three outcomes is new, and Section~\ref{sec:fusiontheory} sets out
the prior claims. What that taxonomy does not provide, having been written for
sensors whose relationship follows from the sensing geometry, is a way to
\emph{determine} which case holds when the sources are a hand-crafted prior
and a learned representation. Their relationship is not evident a priori and,
as Section~\ref{sec:fusion} shows, not predictable from family names: the same
prior is redundant with effective self-supervision and complementary with
augmentation. Our contribution is therefore operational, not taxonomic: fusion adds value exactly when the sources' contributions are
non-overlapping \emph{in feature space}, that overlap is measurable in advance
by an inexpensive linear evaluation of each source alone, and the measurement
predicts the combination. Fusion decisions can be made by measurement, not fashion. On the destructive outcome the grid also locates the damage:
it is carried by the feature term, not the readout, which makes it a
fusion outcome on the same currency logic as the other two, not a separate pathology.

\paragraph{What the attention result does and does not say.}
The persistent, scale-growing ViT deficit is a statement about
\emph{training small-data vision transformers from scratch}, where we
show a fixed spectral target supplies structure that neither 1.28M images
nor the DeiT recipe nor $2\times$-compute SSL supplies at comparable cost. It is
not a claim about web-scale pre-training, where data eventually buys
everything (our own convolutional cells show exactly that limit). The
right reading is deployment-oriented: wherever from-scratch ViTs are
trained on $10^4$--$10^6$ images (medical imaging, remote sensing,
industrial inspection), a $2\%$-overhead structural prior is close to
free accuracy, and the smaller-model-plus-prior configuration can beat
the larger model outright.

\paragraph{On being wrong in public.}
Four of our major pre-registered predictions missed: SSL did not collapse
at 1\% on convolutional backbones (it won everywhere under plain
augmentation); augmentation did not substitute for the prior (it
amplified it); stronger contrastive views did not help SimCLR (they hurt
it); and an early ``gain tracks deficit'' law was falsified outright and
retracted. We report these because each miss redirected the program, and because a benchmark that cannot be wrong measures nothing.

%% ---------------------------------------------------------------------
\section{Limitations and future directions}\label{sec:limits}

\paragraph{The law's status and domain of validity.} The law is an
\emph{empirical} regularity, not a theorem, and it was fitted to nothing: we
report the measured readout curve and use it directly. Its demonstrated domain
is 13 datasets across six visual domains, backbones from 5.7M to 86M
parameters, images from $32$ to $224$\,px and training sets from 500 to
$1.28$M images, all under one supervised recipe with an auxiliary prior
trained from scratch; cells initialized from self-supervised or ImageNet
weights lie outside that scope and are scored separately. Its credibility rests on prediction, not fit, in the three registered episodes of
Section~\ref{sec:lawpredict}. Nothing here licenses extrapolation to web-scale
pre-training or to architectures whose baselines do not train under this
recipe.

\paragraph{Known boundaries.} The account is informative where the readout
term is large and close to uninformative where it has decayed: below the
crossing the sign is right $94\%$ of the time, above it $68\%$, and
\auditUnresolved{} of \auditScope{} in-scope cells cannot test it at all. Read
it as an account of the scarce-data flank. The one non-classification test we
run, semantic segmentation (Section~\ref{sec:dense}), confirms the positive
branch on a different task and metric and cannot test the negative one, for
the reason given there. The second boundary gets its own subsection below,
because it is the one place where the account is not merely weak but wrong in
a patterned way.

\paragraph{Two confounds we can bound but not remove.} \looseness=-1 The first sits inside
the readout term. $G$ is measured with a linear evaluation fit on the full
training split while the cell trained on a fraction of it, so at $1\%$ the
evaluation holds a hundred times the cell's labels; a probe with more labels
reads more of a feature improvement than the cell's own classifier can
realize. Part of the negative readout below the crossing is therefore that asymmetry, not a learning phenomenon, and baseline accuracy is itself a
monotone proxy for the same ratio. The size is measurable: evaluating the
ImageNet64 checkpoints at $25$, $100$ and $250$ labels per class gives $G$ of
$+1.53$, $+2.38$ and $+2.84$ on ViT-tiny against an end-to-end gain of
$+3.24$, so the residual shrinks as the two budgets converge. Holding one
evaluation budget for every cell is what makes cells comparable, which is what
the audit needs, but the readout term should be read as a gap relative to a declared budget, not as a pure property of training. Re-fitting at each
cell's own budget across the crossing would settle how much of the negative
branch survives, and we have not run it. The second confound is the recipe:
it fixes epochs, not steps, so a $1\%$ cell takes a hundredth of the
optimizer steps a $100\%$ cell does and the $\lambda$ schedule is compressed
with them. The envelope across the data axis is thus a joint statement about
data and optimization budget, not data alone, and
Section~\ref{sec:budget:vit} shows the same confound biting a comparator,
where 200 epochs over $2{,}500$ images left SimCLR step-starved, not data-starved.


\begin{figure}[pos=htbp]
\centering
\includegraphics[width=\linewidth]{figs/exceptions.pdf}
\caption{Every resolvable law-scope cell, end-to-end gain against
frozen-feature gain. On the dashed diagonal accuracy realizes the feature gain
exactly. The shaded quadrant is the patterned failure: features improve
measurably while accuracy falls. Only cells whose baseline sits \emph{above}
the crossing are marked, since below it a negative readout is what the account predicts, not a violation of it. Produced by the same audit as
Table~\ref{tab:audit}, so the figure and the exception list cannot disagree.}
\label{fig:exceptions}
\end{figure}

\subsection{Better features, worse accuracy: the one patterned
failure}\label{sec:exceptions}
The decomposition says accuracy tracks features through a readout channel
whose sign is set by baseline height. Above the crossing that channel is
predicted to be neutral, so a cell there should realize its feature gain
almost exactly. Seven resolvable cells do the opposite: the frozen features
measurably improve while end-to-end accuracy falls
(Fig.~\ref{fig:exceptions}).

They do not scatter. Five of the seven are Food-101 and PathMNIST at
${\geq}20\%$ of their data and carry the largest magnitudes
(Table~\ref{tab:exceptions}); the remaining two, a MobileNetV3 EuroSAT cell
and a ResNet-50 CIFAR-100 cell, sit near the threshold at
$G \approx +0.5$--$0.9$ and are small enough that we would not build on them.
The phenomenon is concentrated on two populations at sufficiency rather than
confined to them.

Two things are worth separating. It is \emph{not} the deep-left-flank
behaviour that superficially resembles it: nine further cells, seven on CUB-200
and two on CIFAR-100, also show positive $G$ with $\Delta \leq 0$, but their
baselines sit far \emph{below} the crossing, where the account predicts a
strongly negative readout, so those cells are the account working rather than
failing, and Fig.~\ref{fig:exceptions} excludes them for that reason. And it
is not a measurement artifact on PathMNIST alone, whose linear evaluation we
already flag as a compressed measuring stick; Food-101's cells carry
five times the magnitude and no such caveat.

What is left is a real regime the decomposition has no term for: at
sufficiency, on fine-grained and texture-dominated data, the prior can buy
representation quality that the network then fails to convert. The
overshoot account we use elsewhere, early $\lambda$ shaping costing accuracy
where data is already adequate, predicts the sign of $\Delta$
but not the simultaneous \emph{rise} in $G$, which is what makes these cells informative, not merely inconvenient. Instrumenting the
late-training dynamics of these specific cells is the clearest open problem
the grid poses, and we have deliberately left them in every count in
Section~\ref{sec:lawaudit} rather than scoping them out.

\paragraph{Scope of the instrument.} \looseness=-1 The frozen recipe is, in effect, a
ResNet recipe: architectures that cannot train under it (ConvNeXt, Swin
under SGD) enter only through diagnostic AdamW cells, and their
bistability, itself a finding, limits headline claims. PathMNIST's
shifted test split confounds its high-data cells for every method. Linear
evaluations lose interpretability at 100\% cells (the evaluation ceiling)
and at ImageNet64
are budget-bound by construction. Both ImageNet stages were
subsequently run across $1$--$25\%$ as well as at full data, so the envelope
and not only the right flank is measured at scale; a $1\%$ ImageNet64 cell is
$12{,}820$ images over $1000$ classes, thirteen per class in a label space
five times finer than anything else here. That corner cost us a registered
prediction of an interior peak at $5$--$20\%$, for the reason given in
Section~\ref{sec:scaleenvelope}. One prior family (Gabor and moment energy)
is tested as the knowledge source; the taxonomy predicts, but does not
yet demonstrate, that other structural priors will slot into the same
currency logic.

\paragraph{Future directions.} Three follow directly. (i) \emph{The
exception cluster}: instrument the late-training dynamics of high-data
Food-101 and PathMNIST cells to locate where good features fail to become
accuracy. (ii) \emph{Other currencies}: repeat the fusion matrix with
priors carrying genuinely different structure (depth, symmetry,
frequency-phase) to test whether currency-overlap measured by linear evaluation
predicts stacking in general. (iii) \emph{Attention at deployment
scale}: the model-scale trend ($+13$ rising to $+26$ from ViT-S to ViT-B at fixed
data) invites the obvious next point on the curve, and the
stabilization-plus-structure account suggests trying the prior as a
warmup for large-ViT pre-training proper.

%% ---------------------------------------------------------------------
\section{Conclusion}\label{sec:conclusion}

We opened with one intervention worth $+26$ points, nothing, and $-17$ points
in three settings, and asked what distinguishes them. Sources trade in
currencies: what one is worth is set by which currency is scarce rather than
by how modern it is, and each source's currency is measurable on frozen
features from that source \emph{alone}, before any combination is trained.
Same currency and the second source is redundant however sophisticated;
different currencies and the gains compound; a shaping prior poured into an
initialization that already carries mature features destroys them in
proportion to what they were worth. One decomposition separating feature gain
from readout organizes the grid behind this, holding in sign on
$\auditRate\%$ of testable cells, and the same discipline cost us two accuracy
claims, restated here as cost claims.

The most surprising survivor is the attention result. A bank of oriented
filters that predates deep learning is worth $+13$ points to ViT-S/16 and
$+26$ to ViT-B/16 at $224$\,px, \emph{more} to the larger model, and still
$+3.2$ after 1.28 million images, where our convolutional baseline is neutral
to within $0.04$. Under every recipe we evaluated, neither more parameters nor more data closed
that gap, and $2\%$ of a training run did.
